% =====================================================================
%  Climbing Wall Force-Measurement System
%  User Guide  --  using the dashboard
%
%  Compile with:   latexmk -pdf -synctex=1 user_guide.tex
%  Screenshots live in ./images/user_guide/
% =====================================================================
\documentclass[11pt,a4paper]{article}
\usepackage{climbingwall}
\graphicspath{{images/user_guide/}}
\setdocname{User Guide}
\renewcommand{\shotplacement}{H}  % pin screenshots in reading order (figure-dense)

\begin{document}

\guidetitle{Using the Climbing Wall Force Dashboard}%
  {Getting-started guide --- connecting, recording, and reading the data.}

The climbing wall is fitted with four force sensors --- one under each hand hold
and each foot hold. This web dashboard connects to the sensors, shows the
forces in real time, and lets you record sessions, run a jump-height test, and
compare past recordings. This guide shows you how to start the system and how to
use each feature. It assumes only that you are comfortable opening a terminal and
a web browser; no specialist knowledge is required.

% =====================================================================
\section{Starting the system}
% =====================================================================

The system has three parts that work together: the four \textbf{sensor boards}
(one per hold, plugged into the laptop over USB), a \textbf{backend} (a small
Python server that reads the boards, streams their data, stores recordings and does
the jump calculation) and a \textbf{frontend} (the dashboard you see in the
browser). Set them up in the order below --- the two software parts each in their
own terminal window. To explore in \textbf{demo mode} you can leave the sensor
boards unplugged --- the backend and frontend alone are enough.

\subsection*{1 --- Plug in the sensor boards (live data only)}
The four sensor boards connect to the laptop through a \textbf{USB hub}: each
board's USB cable goes into the hub, and the hub plugs into the laptop with a
single cable. That is all --- there is nothing to log into and nothing to start on
the sensor side. The backend (next step) reads the boards directly over USB.

\begin{infobox}
\textbf{First time on a new computer:} the Phidget sensor driver must be set up
once. On \textbf{macOS}, allow it under \textsf{System Settings $\rightarrow$
Privacy \& Security} (click \textsf{Allow} next to the Phidgets driver). On
\textbf{Windows}, install the Phidget22 drivers from
\texttt{phidgets.com/docs/OS\_-\_Windows}. Without the driver the backend cannot
open the boards (every channel fails with ``device is in use'').
\end{infobox}

Boards may be plugged in or re-plugged at any time, before or after the backend is
started --- they are picked up automatically within a second or two. Only if you
restart the \emph{backend} itself do you need to re-click \textsf{Connect Device}
on the website afterwards.

\subsection*{2 --- Start the backend (Python)}
In a terminal:
\begin{quote}\ttfamily
>> cd backend\\
>> python3 -m venv venv\hfill\textrm{\small(first time only)}\\
>> source venv/bin/activate\hfill\textrm{\small(Windows: venv\textbackslash Scripts\textbackslash activate)}\\
>> pip install -r requirements.txt\hfill\textrm{\small(first time only)}\\
>> python app.py
\end{quote}
You should see \texttt{Running on http://127.0.0.1:5001}. Leave this window open.
There is also a one-step launcher in the \texttt{backend} folder that even does the
first-time setup for you: \texttt{run\_backend.bat} on Windows (double-click it, or
run it from a terminal), \texttt{./run\_backend.sh} on macOS/Linux. Use the one that
matches your computer --- a \texttt{.sh} script will not run on Windows (Windows
asks ``Open with?'' instead) and a \texttt{.bat} will not run on macOS.

\begin{warnbox}
\textbf{Run only one backend at a time.} Each sensor board can be opened by exactly
one program, so a second backend instance --- or the Phidget Control Panel app left
open --- will fight over the boards and neither will get data.
\end{warnbox}

\subsection*{3 --- Start the frontend (Node)}
In a second terminal:
\begin{quote}\ttfamily
>> cd climbing\_wall\_website\\
>> npm install\hfill\textrm{\small(first time only)}\\
>> npm run dev
\end{quote}
This prints a local address such as \texttt{http://localhost:5173}.

\subsection*{4 --- Open the dashboard}
Open the address printed above in a web browser. Any modern desktop browser works
(Chrome, Edge, Firefox, Safari) --- the dashboard receives the sensor data from the
backend, so it no longer needs any special browser feature. Without the sensor
boards it still runs in \textbf{demo mode} with simulated data --- handy for trying
things out.

% =====================================================================
\section{The dashboard at a glance}
% =====================================================================

When the page loads you see the toolbar across the top, a status bar, and four
tabs (Figure~\ref{fig:overall_force}): the three described in this guide, plus a
\textsf{Jump Test (BETA)} tab that is still experimental and not covered here.

\screenshot{overall_force}{0.92}{The dashboard: toolbar (top), status bar, the four
tabs, and the \emph{Overall Force} chart. Here the system is in demo mode and
recording.}

\paragraph{The toolbar.} The controls are grouped into three trays that read left
to right:
\begin{itemize}[leftmargin=1.4em]
  \item \textbf{Session} --- \textsf{Connect Device} opens or closes the live
        sensor stream from the backend (it turns into \textsf{Disconnect} once
        connected), and
        \textsf{Start Recording} / \textsf{Stop Recording} begins or pauses
        collecting data. While paused you also get \textsf{Start New Recording},
        \textsf{Save} and \textsf{Export Data}.
  \item \textbf{Sensor calibration} --- \textsf{Zero Sensors} sets the live zero
        (the ``tare''), \textsf{Calibrate} opens the calibration wizard, and
        \textsf{Profiles} saves or loads named calibrations. All three are covered
        in the separate \emph{Calibration Guide}.
  \item \textbf{View} --- \textsf{Sensor frame / View in world frame} switches the
        coordinate system used for the charts (Section~\ref{sec:world}), and
        \textsf{Chart Settings} adjusts how the charts look
        (Section~\ref{sec:chartsettings}).
\end{itemize}

\paragraph{The status bar.}
A coloured dot reports the device: \textcolor{green!55!black}{green} = connected,
amber = demo mode, red = no device. On the right it shows whether data is currently
being recorded.

% =====================================================================
\section{Recording a session}
% =====================================================================

\begin{enumerate}[leftmargin=1.6em]
  \item Click \textsf{Connect Device} to open the live sensor stream (skip this to
        stay in demo mode). If a warning appears saying only some of the 12 sensor
        channels are attached, check the USB cables of the sensors it names ---
        missing channels read 0.
  \item Click \textsf{Start Recording}. Live data begins to flow into the charts
        and the status bar shows ``Recording data\ldots''.
  \item Climb, hang, or push on the holds. The charts update in real time.
  \item Click \textsf{Stop Recording} to pause. Your data is kept.
  \item Click \textsf{Save} to store the session on the backend (give it a name so
        you can find it later), or \textsf{Export Data} to download it as a
        \texttt{.csv} file you can open in Excel.
  \item Click \textsf{Start New Recording} to clear the buffer and begin a fresh
        session.
\end{enumerate}

\begin{infobox}
The dashboard auto-saves the live buffer to the browser every few seconds, so a
page reload loses at most a few seconds of data. ``Save'' (to the backend) and
``Export'' (to a file) are how you keep a session permanently.
\end{infobox}

% =====================================================================
\section{The chart tabs}
% =====================================================================

\subsection{Overall Force}
The first tab (Figure~\ref{fig:overall_force}) shows the \textbf{total} force on
each of the four sensors as one line each. ``Total'' means the combined strength of
the pull regardless of direction (the length of the force vector,
$\sqrt{X^2+Y^2+Z^2}$). Use it for a quick ``how hard is each limb pulling?''
overview. Because it ignores direction, this chart looks the same in either
coordinate frame.

\subsection{Force by Direction}
The second tab (Figure~\ref{fig:force_by_direction}) breaks each sensor's force
into its three components --- $X$, $Y$ and $Z$ --- on a separate small chart per
sensor. This is where you see, for example, how much a hand is pulling
\emph{down} versus \emph{out of the wall}. The toggle chips above the charts let
you show or hide individual sensors, and each chart can be expanded.

\screenshot{force_by_direction}{0.85}{\emph{Force by Direction}: each sensor's $X$
(purple), $Y$ (slate) and $Z$ (gold) components over time.}

% \subsection{Jump Test (BETA)}
% The third tab (Figure~\ref{fig:jump_test}) estimates how high a climber jumped off
% the wall from the force data. The workflow is deliberately simple:
% \begin{enumerate}[leftmargin=1.6em]
%   \item \textbf{Step 1 --- parameters.} Enter the climber's body weight (kg) and
%         the wall angle (degrees), clicking \textsf{Set} after each.
%   \item \textbf{Step 2 --- run.} Click \textsf{Start Jump Test}, perform the jump,
%         then click \textsf{Record Result}.
% \end{enumerate}
% The result appears as a large number in centimetres. If a body weight was given, a
% \emph{mass-independent score} ($\text{height}/\sqrt[3]{\text{mass}}$) is also shown,
% which makes it fairer to compare climbers of different sizes.

% \screenshot{jump_test}{1.0}{The \emph{Jump Test} tab: enter the parameters, run the
% test, and read the jump height.}

% \begin{infobox}
% Behind the scenes the backend takes the force on the hands and feet during the
% jump, removes the climber's body weight, and integrates the leftover acceleration
% twice (acceleration $\rightarrow$ velocity $\rightarrow$ height) to find the peak.
% You do not need to do any of this --- just enter the weight and angle and press the
% buttons.
% \end{infobox}

\subsection{Recordings}
The \emph{Recordings} tab (Figure~\ref{fig:recordings}) lists your \textbf{five
most recent} saved sessions --- older ones are not shown, but they are not lost
either: every save stays on disk in the backend's \texttt{saved\_recordings}
folder. Click a recording to plot it: you get its \emph{Overall Force} chart and
the per-sensor \emph{Force by Direction} charts. From here you can:
\begin{itemize}[leftmargin=1.4em]
  \item toggle each recording between \textbf{Sensor} and \textbf{World} view (the
        world angle is taken automatically from the tilt stored with the
        recording);
  \item \textsf{Export CSV} to download the session;
  \item \textsf{Compare} two recordings side by side (available once you have at
        least two saved).
\end{itemize}

\screenshot{recordings}{0.8}{The \emph{Recordings} tab showing a saved climbing
session. The clear spike is a hard pull on one hold.}

% =====================================================================
\section{Sensor frame vs.\ world frame}
\label{sec:world}
% =====================================================================

The sensors are bolted to the (tilted) wall, so by default they report forces in
the \textbf{sensor frame}: $X$ runs \emph{along} the wall surface and $Z$ runs
\emph{perpendicular} to it. Often you would rather think in real-world terms ---
straight up versus straight out. The \textbf{world frame} rotates the $X$ and $Z$
components by the wall's tilt angle $\theta$ so that:
\[
  X \rightarrow \text{true vertical}, \qquad Z \rightarrow \text{true horizontal
  (out of the wall)}, \qquad Y \rightarrow \text{unchanged}.
\]

Geometrically the two frames share a single origin and differ only by the tilt:
the world frame is the sensor frame rotated by $\theta$ (Figure~\ref{fig:frames}).

\begin{figure}[H]
\centering
\begin{tikzpicture}[scale=1.05,>=latex,line join=round,
    worldax/.style={->,very thick,accent},
    sensorax/.style={->,very thick,orange!85!black}]
  \def\th{26}
  \coordinate (O) at (0,0);

  % the tilted climbing wall, a soft slab along the sensor-X (up-slope) line
  \draw[line width=8pt,line cap=round,gray!22]
       ($(O)+({1.0*sin(\th)},{-1.0*cos(\th)})$) --
       ($(O)+({-3.0*sin(\th)},{3.0*cos(\th)})$);
  \node[gray!55!black,font=\footnotesize,below right=0pt]
       at ($(O)+({1.0*sin(\th)},{-1.0*cos(\th)})$) {wall surface};

  % a climber on the wall -- fixes which surface "the wall" (sensor X) is
  \begin{scope}[shift={({-2.05*sin(\th)},{2.05*cos(\th)})},rotate=\th,
                scale=0.45,ink,line width=1pt,line cap=round]
    \draw (0,0) -- (0,0.92);                      % spine
    \draw[fill=ink] (0,1.12) circle (0.18);       % head
    \draw (0,0.86) -- (-0.58,1.55);               % left arm reaching to a hold
    \draw (0,0.86) -- (0.52,1.60);                % right arm reaching to a hold
    \draw (0,0) -- (-0.50,-0.78);                 % left leg
    \draw (0,0) -- (0.42,-0.46) -- (0.24,-0.90);  % right leg (bent)
    \fill[gray!55!black] (-0.58,1.55) circle (0.07);  % hold at left hand
    \fill[gray!55!black] (0.52,1.60) circle (0.07);   % hold at right hand
  \end{scope}

  % angle between each world axis and its tilted sensor counterpart
  \draw[gray!60] (0,0.9) arc (90:{90+\th}:0.9);
  \node[gray!45!black,font=\footnotesize]
       at ({1.1*cos(90+\th/2)},{1.1*sin(90+\th/2)}) {$\theta$};
  \draw[gray!60] (0.9,0) arc (0:\th:0.9);
  \node[gray!45!black,font=\footnotesize]
       at ({1.1*cos(\th/2)},{1.1*sin(\th/2)}) {$\theta$};

  % world frame: true vertical + true horizontal
  \draw[worldax] (O) -- (0,2.15) node[accent,above] {$X$};
  \draw[worldax] (O) -- (2.15,0) node[accent,right] {$Z$};

  % sensor frame: along the wall (toward the climber) + perpendicular to the wall
  \draw[sensorax] (O) -- ({-1.45*sin(\th)},{1.45*cos(\th)})
       node[orange!80!black,left] {$X$};
  \draw[sensorax] (O) -- ({1.45*cos(\th)},{1.45*sin(\th)})
       node[orange!80!black,above right] {$Z$};
  \fill (O) circle (1.5pt);
\end{tikzpicture}
\caption{Sensor frame vs.\ world frame, seen side-on, with a climber on the wall to
fix the geometry. The \textcolor{orange!80!black}{sensor frame} is bolted to the
board, so its $X$ runs \emph{along the wall surface} the climber is on and its $Z$
points straight out of it. The \textcolor{accent}{world frame} instead has $X$ truly
vertical and $Z$ truly horizontal --- it is just the sensor frame rotated by the
wall tilt~$\theta$. The $Y$ (sideways) axis points out of the page and is the same in
both frames.}
\label{fig:frames}
\end{figure}

To switch the live view, click \textsf{View in world frame\ldots}, type the
\emph{current} wall tilt, and confirm (Figure~\ref{fig:world_frame}). The view
locks at that angle and can be unlocked at any time. The tilt you type here is a
\textbf{display-only} angle, separate from the calibration angle $\theta$ set in
the calibration wizard: it is merely pre-filled from it, and changing it never
touches the calibration.

\screenshot{world_frame}{0.72}{Opting into the world-frame view: enter and confirm
the current wall tilt $\theta$.}

\begin{infobox}
The world frame only changes \emph{what the charts show}. The stored data always
stays in the raw sensor frame, so switching back and forth never alters or damages
a recording. The \emph{Overall Force} chart is unaffected because total force does
not depend on direction.
\end{infobox}

% =====================================================================
\section{Chart settings}
\label{sec:chartsettings}
% =====================================================================

The \textsf{Chart Settings} button opens a small panel
(Figure~\ref{fig:chart_settings}) with three controls:
\begin{itemize}[leftmargin=1.4em]
  \item \textbf{Time window shown} --- how much recent history to display
        (e.g.\ the last 5 seconds, or everything).
  \item \textbf{Force range (chart top)} --- the maximum value on the vertical
        axis, so you can zoom in or out on the force scale.
  \item \textbf{Auto-adjust scale to data} --- let the chart pick the vertical
        range automatically.
\end{itemize}

\screenshot{chart_settings}{0.78}{\emph{Chart Settings}: time window, force range and
auto-scaling.}

% =====================================================================
\section{Calibration and zeroing}
% =====================================================================

Before the forces shown are accurate in newtons, each board must be
\textbf{calibrated} against known weights --- a one-time procedure done from the
\textsf{Calibrate} button. Separately, you \textbf{zero} the sensors at the start of
each session with the \textsf{Zero Sensors} button, which re-sets the baseline so an
unloaded hold reads zero. Calibration is rare and saved; zeroing is quick, live, and
done often. Both are important enough to have their own document --- please see the
separate \textbf{Calibration Guide} for the full walkthrough (including saved
\textsf{Profiles}) and the reasoning behind it.

% =====================================================================
\section{Troubleshooting}
% =====================================================================

\begin{center}
\begin{tabular}{@{}p{0.42\linewidth}p{0.52\linewidth}@{}}
\toprule
\textbf{Problem} & \textbf{Fix} \\
\midrule
\textsf{Connect Device} fails or falls straight into demo mode &
Make sure the backend is running (\texttt{python app.py} in the \texttt{backend}
terminal) --- the live stream comes through it. \\[4pt]
A warning says only some of the 12 sensor channels are attached &
One or more boards are unplugged. Check the USB cables into the hub for the
sensors named in the warning --- re-plugged boards are picked up automatically
within a second or two. \\[4pt]
Boards are plugged in but no channels attach &
Allow the Phidgets driver once (macOS: \textsf{System Settings $\rightarrow$
Privacy \& Security}) and restart the backend. Also close the Phidget Control
Panel and any second backend instance --- each board can only be opened by one
program at a time. \\[4pt]
``No backend'' / recordings will not save &
Make sure \texttt{python app.py} is still running in the \texttt{backend} terminal. \\[4pt]
\texttt{node} or \texttt{python3} ``command not found'' &
Install Node.js (LTS) from nodejs.org or Python from python.org. \\[4pt]
The charts are empty &
Press \textsf{Start Recording}. With no device connected the system shows demo
data so you can still explore. \\[4pt]
Forces look shifted, or too large or too small &
First \textsf{Zero Sensors} with nothing on the holds. If they are still off, the
sensors probably need calibrating --- see the Calibration Guide. \\
\bottomrule
\end{tabular}
\end{center}

\vfill
\begin{center}\small\itshape\color{muted}
Screenshots were taken from the dashboard running in demo mode with simulated data.
\end{center}

\end{document}
